% Growth-twin manuscript (Path A), IEEE journal format.
% Every number is a macro from ../tables/twin_values.tex or a table from ../tables/twin_tab_*.tex,
% generated by scripts/build_twin_paper.py from runs. Do not type numbers here (results guard).
\documentclass[journal]{IEEEtran}
\usepackage[T1]{fontenc}
\usepackage{amsmath,amssymb}
\usepackage{graphicx}
\usepackage{booktabs}
\usepackage{xcolor}
\usepackage{cite}
\usepackage{url}
\usepackage[hidelinks]{hyperref}
\graphicspath{{../figures/}}
\input{../tables/twin_values.tex}
\providecommand{\PENDINGFIG}{\textcolor{red}{[figure pending]}}

\begin{document}

\title{A Sequentially Updated, Uncertainty-Calibrated Growth Twin for Maturity Monitoring from Sparse Measurements in Youth Athletes}

% TODO(human): affiliations, corresponding author, funding and ethics statements.
\author{Anirudhha~Veeranagaiah~M and Shreyas~J%
\thanks{Code and run records: project repository; every number in this paper is generated from a run.}}

\markboth{Manuscript draft}{Sequential growth twin for maturity monitoring}

\maketitle

\begin{abstract}
Youth sport programmes estimate biological maturity to group players, manage training load and
anticipate growth-related injury, usually with a maturity-offset equation applied to one
measurement. These equations carry no uncertainty, do not learn from repeated measurements, and
are systematically wrong in early and late maturers. We present a growth twin: a per-athlete
Bayesian growth-curve state, learned against a reference cohort and updated after every new
measurement, that estimates age at peak height velocity (PHV) with conformally calibrated
intervals, gives the probability that the athlete is in the growth spurt now, and forecasts the
next measurement. In five-fold validation by child on the real Berkeley Child Guidance Study
(\TwChildren{} children, \TwEvaluations{} visit-by-visit evaluations), the twin's error in age at
PHV was \TwMAE{}~years against \TwMooreMAE{} for the Moore equation (difference \TwDiffMoore{};
95\,\% CI \TwDiffMooreCI{}); the two did not differ at the first visit, and the difference grew to
\TwDiffMooreLast{}~years by the sixth. The advantage was concentrated in early and late maturers
(\TwEarlyDiffMoore{} and \TwLateDiffMoore{}~years against Moore), where equations are known to
fail. Classical SITAR, refitted at every visit, was as accurate overall (difference
\TwDiffSitar{}; \TwDiffSitarCI{}), less accurate in early maturers and more accurate in average
maturers; gradient boosting likewise matched the twin overall (\TwDiffML{}; \TwDiffMLCI{}). The twin's
distinctive contribution is its per-child uncertainty. Its calibrated intervals covered \TwCovConf{}
at nominal \SpecInterval{} and adapted to the child: its error rose from \ErrNarrow{} to
\ErrWide{}~years from its narrowest to its widest third of intervals, and in that widest third
its intervals covered \CovTwinWide{} against \CovSitarWide{} for intervals of constant width
calibrated in the same way around SITAR. The growth-spurt probability was well
calibrated (slope \CircaSlope{}) and discriminated better than the equation (AUC \TwCircaAUC{}
against \TwCircaAUCMoore{}).
\end{abstract}

\begin{IEEEkeywords}
Digital twin, growth curve, peak height velocity, maturity offset, Bayesian hierarchical model,
conformal prediction, youth sport, bio-banding.
\end{IEEEkeywords}

\section{Introduction}
\IEEEPARstart{B}{iological} maturity differs by several years between children of the same age,
and youth sport uses it to group players by maturity (bio-banding) \cite{cumming2017biobanding},
to adjust training load, and to anticipate growth-related injuries that cluster around the
adolescent growth spurt \cite{hall2025,oxendale2025}. In practice maturity is estimated from one
measurement with a regression equation for maturity offset, the years from peak height velocity
(PHV) \cite{mirwald2002,moore2015}. Independent validations show that these equations are
reasonable only for average-maturing children near PHV: they place PHV too late in early maturers
and too early in late maturers, and an individual's predictions vary considerably from one
measurement to the next \cite{koziel2018modified,malina2021soccer}. They also give no statement of
uncertainty and cannot use repeated measurements.

Growth-curve models such as SITAR \cite{cole2010sitar}, including a Bayesian implementation
\cite{bsitar}, describe each child's curve as a shifted and stretched version of a shared curve and
estimate age at PHV well from long series. They are fitted retrospectively to complete records,
and fits to sparse, irregular data can show convergence or boundary problems
\cite{carvalho2025sitar}; a new child's estimate requires refitting the cohort model, and no
interval is given for it. Programmes, by contrast, see a few measurements at a time and need an
answer, with its uncertainty, at each visit. We use the term digital twin for an individual model of this kind, updated as
data arrive and reporting its uncertainty; our literature search found no such model for
adolescent growth.

We call our approach a growth twin: a per-athlete model state that is learned against a reference
cohort, updated after every new measurement, reports calibrated uncertainty, forecasts the next
measurement, and feeds a decision (whether the athlete is in the growth spurt). Contributions:
(i) a marginalised Bayesian SITAR-type twin that is exact for one to a few measurements, with a
timing prior whose spread grows with distance from PHV; (ii) cross-fold conformal calibration of
its intervals; (iii) an out-of-sample, visit-by-visit evaluation on real longitudinal data against
the equations used in practice, classical SITAR refitted at every visit, and a machine-learning
baseline, with child-level confidence intervals; (iv) an analysis by maturity timing, the
equations' known failure mode, and of whether calibrated intervals hold for each kind of child;
and (v) robustness to the measurement conditions of real programmes and an external check.

\section{Methods}

\subsection{Growth model}
For athlete $i$ with stature $y_{ij}$ at age $t_{ij}$,
\begin{equation}
y_{ij} = \mu + \alpha_i + \delta_i (t_{ij}-\bar t_i) + h\big((t_{ij}-\beta_i)e^{\gamma_i}\big) + \varepsilon_{ij},
\end{equation}
where $h$ is a natural cubic spline shared by athletes of one sex, $\alpha_i$ shifts size,
$\beta_i$ timing, $\gamma_i$ tempo, $\delta_i$ absorbs a mismatch in velocity, and
$\varepsilon_{ij}$ is measurement error. Size and velocity effects are integrated analytically and
timing and tempo numerically on grids, so athlete effects are drawn exactly given the shared
parameters, and two-valued timing posteriors from short series are represented faithfully. Age at
PHV is the age of the highest interior peak of the athlete's velocity curve.

\emph{Distance-dependent timing prior.} The prior for $\beta_i$ is centred on the Moore estimate
\cite{moore2015} with standard deviation $\sigma_\beta(1+\kappa\,|m_i|)$, where $m_i$ is the Moore
maturity offset at the visit used for the prior and $\kappa$ is learned, encoding the known loss of
accuracy of the equations far from PHV \cite{koziel2018modified}.

\subsection{The twin: reference cohort and sequential updating}
The shared parameters are learned from a reference cohort's full records. Each reference child's
prior is computed from one random season of its own record, as it would be at deployment, so the
learned prior spread reflects real prior error. A new athlete's twin is the posterior of its
effects given the shared parameters and the measurements so far; each new visit redraws it. The
shared curve is fitted at its posterior mode, as SITAR fits by maximum likelihood; the athlete's
state remains fully Bayesian given it. One fold was refitted with full MCMC (NUTS) as a check.

\subsection{Calibration and decision output}
Raw posterior intervals are rescaled by cross-fold split conformal calibration
\cite{lei2018conformal,angelopoulos2021gentle}: the score is the error in units of the athlete's
posterior standard deviation, and for athletes in fold $k$ the scaling quantile is taken from the
held-out athletes of the other folds of the same sex with the same number of visits, so no athlete
calibrates its own interval. The same procedure calibrates the \SpecHorizon{}-month stature
forecast. The decision output is the posterior probability that the athlete is within
\SpecCircaWidth{}~year of PHV now.

\subsection{Data}
The Berkeley Child Guidance Study \cite{tuddenham1954}, distributed with the \texttt{sitar}
package \cite{sitarpkg}, follows \TwChildren{} children born in 1928--29 with stature, stem
(sitting) height and mass every six months from 8 to 21 years; it was not used to develop the
Moore equations. Each child's reference age at PHV was obtained from the full record by a
Preece--Baines model 1 fit \cite{preece1978} checked against a smoothing spline. The Chard Growth
Study (12 girls measured six-monthly from 8 to 16 years with age at menarche, no sitting height;
also in \texttt{sitar}) served as an external check.

\subsection{Evaluation and statistical analysis}
Children were split into \TwFolds{} folds by child, separately by sex. For each fold the reference
curve was learned from the other folds; each held-out child entered at a random measurement age
between 8 and 15 years and the twin was updated after each of up to \TwMaxVisits{} six-monthly
visits. After each update we recorded age at PHV (against the Preece--Baines reference, against the
same model given the child's whole record, which removes differences between curve families, and,
as a sensitivity analysis, against the smoothing-spline reference), interval coverage, the
growth-spurt probability, and the forecast of the next measurement.

\emph{Comparators} used the same visits: Moore and, for boys, Mirwald, taken at the visit each
places closest to PHV; extrapolation of the last observed velocity for the forecast; and a
gradient-boosting regressor trained, within each fold, on sequences from the training children
(features: visits seen, age and measurements at the latest visit, observed stature and
sitting-height velocities, years observed, Moore estimate; target: the Preece--Baines reference)
and given the same cross-fold conformal intervals. Classical SITAR \cite{cole2010sitar} was fitted
by maximum likelihood with the \texttt{sitar} package \cite{sitarpkg} (natural cubic spline with
\SitarDF{} degrees of freedom, random size, timing and intensity), as a practitioner would use it:
after each visit it was refitted to the fold's training children's full records together with the
held-out child's visits so far, and the child's age at PHV was the age of the peak of its fitted
velocity curve. As SITAR gives no interval for a new child's age at PHV, its estimate, too, was given
the cross-fold conformal interval.

\emph{Conditional coverage.} Conformal calibration guarantees coverage on average, not for every
kind of child. We therefore report calibrated coverage within maturity-timing groups and within
thirds of the twin's interval width (formed within each sex and number of visits), where intervals
that adapt to the child should hold their coverage and intervals of constant width should not.

\emph{Inference.} Differences are reported with 95\,\% percentile intervals from \TwBootN{}
bootstrap resamples of children, since a child contributes up to \TwMaxVisits{} correlated
evaluations. Maturity timing was classified as early or late when the reference age at PHV lay more
than \SpecTimingSD{}~SD from the mean of the same sex. The growth-spurt probability was assessed by
AUC, Brier score and logistic calibration intercept and slope.

\emph{Robustness} was tested by degrading only the held-out children's data: yearly instead of
six-monthly visits, added measurement error, and no sitting height (Moore's stature-only equation
for boys). A leg-length extension, in which leg length shares each child's timing with an
individual leg-to-stature lag, was evaluated in the same way.

\section{Results}

\subsection{Accuracy as visits accumulate}
Over \TwEvaluations{} evaluations the twin's mean absolute error in age at PHV was \TwMAE{}~years
against \TwMooreMAE{} for Moore and \TwMirwaldMAE{} for Mirwald (Tables~\ref{tab:visits}
and~\ref{tab:comparisons}, Fig.~\ref{fig:sequential}). At the first visit the twin and Moore did not
differ (\TwDiffMooreFirst{}; 95\,\% CI \TwDiffMooreFirstCI{}); the twin's error then fell to
\TwMAELast{} after \TwMaxVisits{} visits, while the equation, which can only choose which single
measurement to use, improved from \TwMooreFirst{} to \TwMooreLast{}, so that the difference reached
\TwDiffMooreLast{} (\TwDiffMooreLastCI{}). Overall the difference was \TwDiffMoore{}
(\TwDiffMooreCI{}); against the full-information reference \TwDiffMooreFull{}
(\TwDiffMooreFullCI{}). In boys the twin was not significantly more accurate than Mirwald
(\TwDiffMirwald{}; \TwDiffMirwaldCI{}). The twin's bias was \TwBiasPB{}~years against
Preece--Baines but \TwBiasFull{} against its own full-record estimate, so the former is a difference
between curve families; against the smoothing-spline reference the errors were \SplineTwin{} for
the twin and \SplineMoore{} for Moore.

\begin{table*}[!t]
\caption{Age at PHV by number of visits seen by the twin (real Berkeley data, held-out children).
Mean absolute error in years; coverage of raw and conformally calibrated nominal \SpecInterval{}
intervals; mean calibrated width in years.}
\label{tab:visits}
\centering
\input{../tables/twin_tab_visits.tex}
\end{table*}

\begin{table}[!t]
\caption{Paired differences (twin minus comparator; negative favours the twin) with child-level
bootstrap 95\,\% intervals.}
\label{tab:comparisons}
\centering
\resizebox{\columnwidth}{!}{\input{../tables/twin_tab_comparisons.tex}}
\end{table}

\begin{figure*}[!t]
\centering
\IfFileExists{../figures/twin_fig_sequential.png}{\includegraphics[width=\textwidth]{twin_fig_sequential.png}}{\fbox{\PENDINGFIG}}
\caption{Visit-by-visit performance on real held-out children: error in age at PHV (twin and Moore),
raw interval coverage, and six-month stature forecast error (twin and velocity extrapolation).}
\label{fig:sequential}
\end{figure*}

\subsection{Early and late maturers}
Table~\ref{tab:timing} shows where the gains arise. Moore's error was large and directional in
early and late maturers (late estimates in early maturers, early estimates in late maturers), as
reported in independent samples \cite{koziel2018modified}. The twin reduced the error by
\TwEarlyDiffMoore{}~years (\TwEarlyDiffMooreCI{}) in the \TwEarlyChildren{} early maturers and by
\TwLateDiffMoore{} (\TwLateDiffMooreCI{}) in the \TwLateChildren{} late maturers, while in the
\TwAverageChildren{} average maturers the two did not differ significantly (\TwAverageDiffMoore{};
\TwAverageDiffMooreCI{}). The twin's own bias in these groups was smaller than Moore's but not zero.

\begin{table}[!t]
\caption{Mean absolute error (mean bias) in years by maturity timing; early and late beyond
\SpecTimingSD{}~SD of the sex-specific mean reference age at PHV.}
\label{tab:timing}
\centering
\resizebox{\columnwidth}{!}{\input{../tables/twin_tab_timing.tex}}
\end{table}

\subsection{Calibrated uncertainty and the growth-spurt probability}
Raw nominal \SpecInterval{} intervals covered \TwCovRaw{} of cases (mean width \TwWidthRaw{}~years);
after conformal calibration they covered \TwCovConf{} (width \TwWidthConf{}), with coverage close to
nominal for each sex and number of visits (Table~\ref{tab:visits}) but lower in late maturers
(Table~\ref{tab:timing}). The learned timing prior widened by \TwSlope{} of its standard deviation
per year of distance from PHV. \TwCircaPrev{} of evaluations were within \SpecCircaWidth{}~year of
PHV; the twin's probability discriminated these with an AUC of \TwCircaAUC{} against
\TwCircaAUCMoore{} for the Moore offset (difference \TwDiffAUC{}; \TwDiffAUCCI{}), and was well
calibrated (intercept \CircaIntercept{}, slope \CircaSlope{}, Brier score \CircaBrier{}). At a
probability of one half its sensitivity and specificity were \TwCircaSensTwin{} and
\TwCircaSpecTwin{}, against \TwCircaSensMoore{} and \TwCircaSpecMoore{} for an offset within one year.

\subsection{Machine-learning baseline}
Gradient boosting trained on the reference children reached an error of \MlMAE{}~years, matching the
twin overall (difference \TwDiffML{}; \TwDiffMLCI{}) and at early visits (first visit \TwDiffMLFirst{};
\TwDiffMLFirstCI{}), with calibrated coverage of \MlCov{} at a similar width (\MlWidth{}~years) and a
growth-spurt AUC of \MlAUC{} (difference from the twin \TwDiffMLAUC{}; \TwDiffMLAUCCI{}). The twin was
more accurate after six visits (\TwDiffMLLast{}; \TwDiffMLLastCI{}) and in early maturers
(\TwEarlyDiffML{}; \TwEarlyDiffMLCI{}), and against the spline reference (\SplineTwin{} against
\SplineML{}); in average and late maturers the two did not differ significantly.

\subsection{Classical SITAR}
All \SitarRows{} SITAR refits converged, helped by fitting each sparse child
together with the reference children's full records. SITAR's error was \SitarMAE{}~years against
\SitarTwinSame{} for the twin on the same evaluations (difference \TwDiffSitar{}; \TwDiffSitarCI{}),
with no significant difference at the first visit (\TwDiffSitarFirst{}; \TwDiffSitarFirstCI{}) or
the sixth (\TwDiffSitarLast{}; \TwDiffSitarLastCI{}), nor in growth-spurt discrimination (AUC
\SitarAUC{}; difference \TwDiffSitarAUC{}; \TwDiffSitarAUCCI{}). The twin was more accurate in
early maturers (\TwEarlyDiffSitar{}; \TwEarlyDiffSitarCI{}) and less accurate in average maturers
(\TwAverageDiffSitar{}; \TwAverageDiffSitarCI{}); in late maturers they did not differ
(\TwLateDiffSitar{}; \TwLateDiffSitarCI{}). Against the smoothing-spline reference SITAR's error
was \SplineSitar{} and the twin's \SplineTwinSitarRows{}. With the conformal interval SITAR's
coverage was \SitarCov{} at a width of \SitarWidth{}~years.

\subsection{Do calibrated intervals hold for each child?}
Average coverage was close to nominal for all three methods, but not within subgroups
(Table~\ref{tab:coverage}). The twin's posterior standard deviation tracked its error (Spearman
\AdaptRho{}): from its narrowest to its widest third of intervals the twin's error rose from
\ErrNarrow{} to \ErrWide{}~years and its coverage was \CovTwinNarrow{}, \CovTwinMiddle{} and
\CovTwinWide{}, somewhat below nominal in the narrowest third. SITAR's constant-width intervals
varied more and in the opposite direction: they over-covered the easy cases (\CovSitarNarrow{};
difference from the twin \CovNarrowDiffSitar{}, \CovNarrowDiffSitarCI{}) and under-covered the hard
ones (\CovSitarWide{}; \CovWideDiffSitar{}, \CovWideDiffSitarCI{}). In early maturers the twin
covered \CovTwinEarly{} against \CovSitarEarly{} for SITAR and \CovMlEarly{} for boosting, a
difference from SITAR (\CovEarlyDiffSitar{}; \CovEarlyDiffSitarCI{}) that is not significant in
\TwEarlyChildren{} children. In late
maturers all three under-covered (twin \CovTwinLate{}, SITAR \CovSitarLate{}, boosting
\CovMlLate{}; difference from SITAR \CovLateDiffSitar{}, \CovLateDiffSitarCI{}), and in average
maturers all three over-covered (twin \CovTwinAverage{}, SITAR \CovSitarAverage{}).

\begin{table}[!t]
\caption{Coverage of calibrated nominal \SpecInterval{} intervals within subgroups (real data,
held-out children). Width thirds are formed from the twin's interval within each sex and number of
visits; MAE and width in years.}
\label{tab:coverage}
\centering
\resizebox{\columnwidth}{!}{\input{../tables/twin_tab_coverage.tex}}
\end{table}

\subsection{Forecasting the next measurement}
The twin's \SpecHorizon{}-month forecast error was \TwForecastMAESame{}~cm against \TwNaiveMAE{} for
extrapolating the last observed velocity on the same visits. Calibrated forecast intervals covered
\TwForecastCovConf{} of measurements (raw \TwForecastCovRaw{}) with a width of
\TwForecastWidthConf{}~cm.

\subsection{Robustness and model variants}
With yearly visits the error was \TwYearlyMAE{}~years (Moore \TwYearlyMoore{}), lower than with
six-monthly visits because the same number of yearly visits spans twice as much growth; with one
centimetre of added measurement error \TwNoiseMAE{} (Moore \TwNoiseMoore{}); without sitting height
\TwNosittingMAE{} (Moore stature-only \TwNosittingMoore{}). Calibrated coverage stayed at
\TwYearlyCov{}, \TwNoiseCov{} and \TwNosittingCov{} (Table~\ref{tab:variants}). Adding leg length did
not help: \TwLegSharedMAE{}~years with shared timing and \TwLegLagMAE{} with an individual lag,
against \TwMainMAE{} for stature alone.

\begin{table}[!t]
\caption{Robustness (held-out data degraded) and leg-length variants (real data).}
\label{tab:variants}
\centering
\resizebox{\columnwidth}{!}{\input{../tables/twin_tab_variants.tex}}
\end{table}

\subsection{Full MCMC check}
Refitting the shared curve of one fold by NUTS gave an error of \TwNutsMAE{}~years and raw coverage
\TwNutsCov{}, against \TwMapFoldMAE{} and \TwMapFoldCov{} with the posterior mode on the same
children (largest R-hat \TwNutsRhat{}, smallest bulk effective sample size \TwNutsESS{}).
The two fits agree closely. With the short chains that were feasible here the sampler had not fully
converged by the usual R-hat criterion, so this is a consistency check of the posterior-mode
approximation rather than a full-MCMC evaluation.

\subsection{External check}
Trained on all Berkeley girls and applied visit by visit to the \ChGirls{} Chard girls, the twin's
final estimate differed from the full-record age at PHV by \ChTwinMAE{}~years against \ChMooreMAE{}
for Moore (\ChRef{} girls with a full-record reference). In the \ChMen{} girls with observed
menarche, the twin's age at PHV correlated \ChCorrTwin{} with age at menarche (Moore
\ChCorrMoore{}), and menarche followed the twin's PHV by \ChLagMean{}~years (SD \ChLagSD{}; Moore
\ChMooreLagSD{}). These samples are too small for inference and are reported as supporting evidence.

\section{Discussion}
The equations used in youth sport are adequate for average maturers near PHV and fail where it
matters most: in early and late maturers, the children bio-banding is designed for and in whom the
timing of growth-related injury risk is least predictable. The growth twin's advantage is
concentrated exactly there, and it grows as visits accumulate, which is how programmes collect
data. It also does what an equation cannot: it states its uncertainty, and once conformally
calibrated its intervals and its growth-spurt probability mean what they say.

The comparisons with classical SITAR and gradient boosting sharpen the claim. Both were as accurate
on average, so the twin's value is not raw accuracy: once repeated measurements are modelled with a
shared growth curve, by maximum likelihood or Bayesian updating, point estimates are similar, and
SITAR was somewhat better in average maturers. The twin's distinctive contribution is uncertainty
that is specific to the child. Conformal calibration makes any estimator's intervals right on
average, but intervals of constant width are too wide for the children who are easy to place and
too narrow for those who are not. The twin's posterior widths separated these children, so its
calibrated coverage varied much less across them and was significantly higher for the hardest
cases; the same direction in early maturers, a group of particular concern for bio-banding and
growth-related injury, needs a larger sample to confirm. Child-specific uncertainty is what makes
an estimate usable for decisions about an individual. The twin also updates a new child without
refitting the cohort model, and against boosting it was more accurate after six visits and in early
maturers and against the spline reference.

Several findings limit the claims. The twin gave no gain over Moore at the first visit, was not
significantly better than Mirwald in boys, and was not more accurate than classical SITAR overall.
Calibrated coverage was below nominal in late maturers for all methods, and the twin's
intervals were somewhat too narrow for the children it placed most confidently. Its
six-month forecast was only on a par with velocity extrapolation, and leg length, despite the
known earlier peak of the legs, did not improve the estimate because the leg-to-stature lag varies
between children.

\emph{Limitations.} The development data come from a historical North American cohort; the twin has
not been tested in athletes or in Indian children, where growth timing differs, and the external
check is small. Classical SITAR was benchmarked on stature only and with one spline basis; we did
not benchmark Bayesian SITAR \cite{bsitar}, which could provide child-level intervals of its own,
nor the maturity-ratio equation of Fransen et al.; the reference age at PHV is itself a model
estimate; and the shared curve was fitted at its posterior mode, with full MCMC checked on one
fold.

\section{Conclusion}
A sequentially updated growth twin can replace one-off maturity equations in youth sport programmes
where it matters most: in early and late maturers, and as measurements accumulate. It is as accurate
as classical SITAR refitted at every visit, and adds what SITAR does not give: calibrated
uncertainty that is specific to each child, and a probability that the athlete is in the growth
spurt now.

\section*{Data and Code Availability}
Berkeley and Chard data: \texttt{sitar} R package (CRAN). Code, configuration and run records are in
the project repository; every number in this paper is generated from a run and traceable to its
configuration, seed and commit. A template for applying the twin to a programme's own records is
included.

\bibliographystyle{IEEEtran}
\bibliography{references}

\end{document}
